%=============================================================================
% 丘成桐中学科学奖（中国内地赛区）研究报告 LaTeX 模板
% ---------------------------------------------------------------------------
% 结构以官方《研究报告模板》（yau-awards.com 文件下载页，2025-04-14 版，4 页）为纲：
%   P1 封面页 → P2 起 题目/作者/摘要/关键词/目录/正文 → 另起一页 参考文献 → 致谢页
% 官方规则依据（引用处在各节注释中标明）：
%   - 参赛规则 https://www.yau-awards.com/page-rule.html  四（四）研究报告格式
%   - AI 使用规则 https://www.yau-awards.com/show-86-59.html（2026-07-08）
%   - 六科评审标准 https://www.yau-awards.com/page-criteria.html（报告八要素）
% 注意：官方模板致谢页写“1-2 页（500-1000 字）”，但 2026 参赛规则页写 500-1500 字
% 且必须含 AI 使用说明——以规则页为准（本模板按 500-1500 字执行）。
%
% 编译：xelatex（需编译两遍以生成目录）。字体用 macOS 自带 Songti SC / PingFang SC；
% Windows/Linux 用户请把下方字体改为 SimSun / SimHei 等本机可用字体。
% 所有〈尖括号〉内容为占位说明，使用前须整体替换并删除说明文字。
%=============================================================================
\documentclass[12pt, a4paper]{article}

\usepackage[margin=2.5cm]{geometry}
\usepackage{xeCJK}
\setCJKmainfont[BoldFont={PingFang SC Semibold}]{Songti SC}
\setCJKsansfont{PingFang SC}
\setmainfont{Times New Roman}
\xeCJKDeclareCharClass{CJK}{"2460 -> "2473}% ①②③④ 等带圈数字交给中文字体
\usepackage{amsmath, amssymb}
\usepackage{graphicx}
\usepackage{booktabs}
\usepackage{caption}
\usepackage[hidelinks]{hyperref}
\usepackage{indentfirst}
\setlength{\parindent}{2em}
\linespread{1.5}

% 图表题注中文化
\renewcommand{\figurename}{图}
\renewcommand{\tablename}{表}
\renewcommand{\contentsname}{目录}
\renewcommand{\refname}{参考文献}
\renewcommand{\abstractname}{摘要}

% 封面填空线（占位说明用小字号，避免溢出下划线）
\newcommand{\coverline}[1]{\underline{\makebox[8.6cm][c]{\normalsize #1}}}

\begin{document}

%=============================================================================
% 封面页 —— 官方模板第 1 页
% 依据：page-rule.html 四（四）2(2)a)“封面页：姓名、省份/州、国别、指导老师姓名、
% 报告标题”；官方模板 PDF 第 1 页实际含以下 7 行（五项规则字段 + 中学 + 指导老师单位），
% 以官方模板 PDF 为准，逐项保留。
%=============================================================================
\begin{titlepage}
  \vspace*{3cm}
  \begin{center}
  {\Large
  \begin{tabular}{@{}rl@{}}
    参赛学生姓名： & \coverline{〈队员姓名，1--3 人〉} \\[1.6em]
    中\quad\quad\quad 学： & \coverline{〈学校全称，须同一中学〉} \\[1.6em]
    省\quad\quad\quad 份： & \coverline{〈省/直辖市；海外写州〉} \\[1.6em]
    国家/地区： & \coverline{〈如：中国〉} \\[1.6em]
    指导老师姓名： & \coverline{〈1--2 位〉} \\[1.6em]
    指导老师单位： & \coverline{〈与盖章单位一致〉} \\[1.6em]
    论文题目： & \coverline{〈与报名系统所填一致〉} \\
  \end{tabular}
  }
  \end{center}
  \vfill
\end{titlepage}

%=============================================================================
% 第二页起：题目、作者、摘要、关键词 —— 官方模板第 2 页
% 依据：page-rule.html 四（四）2(2)b)“第二页起：题目、作者、摘要、关键词、目录、正文”。
%=============================================================================
\begin{center}
  {\LARGE\bfseries 〈论文题目：与封面、报名系统完全一致〉}\\[1.5em]
  {\large 〈作者姓名，顿号分隔，顺序与封面一致〉}\\[0.5em]
  {\normalsize 〈学校名称〉}
\end{center}

\begin{abstract}
〈摘要：研究问题、方法、主要结果、结论各 1--3 句，写实际做出的东西，不写空话。
数学奖注意：数学评审标准要求摘要不少于 1 页，且含背景综述与引文，
明确区分“学界已知”与“本文原创”（来源：page-criteria.html 数学奖）。
其他学科摘要通常 200--500 字。〉
\end{abstract}

\noindent\textbf{关键词：}〈3--6 个，分号分隔〉

%=============================================================================
% 目录 —— 官方模板第 2 页要素之五
%=============================================================================
\clearpage
\tableofcontents
\clearpage

%=============================================================================
% 论文正文 —— 官方模板第 2 页要素之六
% 章节骨架按六科评审标准的“报告八要素”组织（page-criteria.html：题目、摘要、
% 背景与现状、材料与方法、结果、讨论、参考文献、每位作者贡献——化学/生物/计算机/
% 经济金融建模明文列出；数学、物理表述为“报告规范性”，内容相同）。
% 题目、摘要已在上文；“每位作者贡献”按官方模板归入致谢页；其余四要素即下列四节。
% 各学科可按需改节名（如数学用“预备知识/主要结果/证明”），但要素不可缺。
%=============================================================================
\section{引言（背景与现状）}
〈研究背景：这个问题为什么值得做；前人做到了哪里（引文标注 \cite{example-ref-1}）；
本文的切入点与贡献。评审标准共同项要求“能区分学界已知与本项目原创”——
在引言末尾用一段明确写出本文的原创点与全文结构安排。〉

\section{材料与方法（模型与理论）}
〈实验类：材料、仪器（型号）、实验设计、控制变量、步骤；须可复核（物理奖评审标准
将“实验过程”列入科学性与严密性）。理论/建模类：模型假设、推导、算法。
计算机奖须给出问题定义与算法实现细节；理论类须完备证明（page-criteria.html 计算机奖）。〉

\subsection{〈子节示例：模型假设〉}
〈……〉

\section{结果}
〈实际得到的数据、定理、系统与图表。化学、生物奖明文规定“仅提供一个研究思路的论文
不得参与评奖，必须要有实质性的研究成果”（page-criteria.html）——本节不可空。
图表规范示例见图~\ref{fig:example} 与表~\ref{tab:example}。〉

\begin{figure}[htbp]
  \centering
  \fbox{\parbox[c][4cm][c]{0.7\textwidth}{\centering 〈图片占位：替换为
  \texttt{\textbackslash includegraphics} 〉}}
  \caption{〈图注：一句话说明图的内容与结论〉}
  \label{fig:example}
\end{figure}

\begin{table}[htbp]
  \centering
  \caption{〈表注：一句话说明表的内容〉}
  \label{tab:example}
  \begin{tabular}{lcc}
    \toprule
    〈变量〉 & 〈条件 1〉 & 〈条件 2〉 \\
    \midrule
    〈行〉 & 〈值〉 & 〈值〉 \\
    \bottomrule
  \end{tabular}
\end{table}

\section{讨论（结论与展望）}
\subsection{结论与意义}
〈结果说明了什么；与前人工作的差异点（数学奖标准：“照搬已有方法竞争力不足”，
须写明与前人方法的差异）；局限性（经济金融建模奖：纯理论推演无实证的要主动说明局限）。〉
\subsection{未来展望}
〈下一步可做什么。〉

%=============================================================================
% 参考文献 —— 另起一页（官方模板第 3 页）
% 依据：page-rule.html 四（四）2(2)c)“另起一页：参考文献”。
% 简单版用 thebibliography；条目多可换 biblatex（\usepackage[style=numeric]{biblatex}）。
%=============================================================================
\clearpage
\begin{thebibliography}{99}
\bibitem{example-ref-1} 〈作者. 题名. 期刊/出版社, 年份, 卷(期): 页码.〉
\bibitem{example-ref-2} 〈网络资源注明 URL 与访问日期；不得出现 AI 生成的虚假文献
（show-86-59.html 禁止“生成虚假文献”，违规取消参赛资格）。〉
\end{thebibliography}

%=============================================================================
% 致谢页 —— 另起一页（官方模板第 4 页）
% 依据：page-rule.html 四（四）2(2)d) + 致谢页条款：1-2 页，500-1500 字
% （官方模板旧写 500-1000 字，以规则页 500-1500 为准）。四项内容缺一不可：
%   ① 整体研究过程逐环节分工：选题来源、数据获取、分析数据及计算、实验设计及实施、
%      撰写论文——逐人具体到环节，并说明遇到的困难及解决经过（评审标准共同项
%      “团队合作/每位作者贡献”即对应此处；计算机奖强调“避免搭车参赛”）；
%   ② 指导老师与参赛学生的关系、在论文写作过程中所起的作用、指导是否有偿；
%   ③ 他人协助完成的研究成果（化学奖标准：大型仪器等他人帮助须注明帮助者姓名与
%      具体内容）；
%   ④ AI 使用情况（show-86-59.html，2026-07-08）：工具名称及版本、具体使用环节及
%      用途、使用时间及频率；使用须事先获指导教师许可；聊天记录及辅助材料上传报名
%      系统“其他材料”栏。未使用 AI 也要写明“未使用”。
% 完整可读的写法示例见同目录《致谢页范例.md / .docx》。
%=============================================================================
\clearpage
\section*{致谢}
\addcontentsline{toc}{section}{致谢}

〈开头 1--2 句：研究背景与感谢总起（可含选题来源的故事）。〉

〈① 分工说明：逐环节、逐人。示例句式——“选题来源于……；实验装置由〈甲〉设计搭建，
〈乙〉负责数据采集与误差分析；理论推导由〈甲〉完成，〈乙〉复核；论文初稿由〈乙〉执笔，
两人共同修订校对。在〈某环节〉我们遇到〈困难〉，通过〈做法〉解决。”〉

〈② 导师说明：“指导老师〈姓名〉系〈本校物理教师 / ×× 大学 ×× 学院研究员〉，
在〈选题讨论 / 理论把关 / 论文修改意见〉环节给予指导，具体的〈实验、计算、写作〉
均由参赛学生独立完成；指导为无偿。”——有偿必须如实写明；关系表述须与
学术诚信声明、指导老师信息表一致。〉

〈③ 他人协助：“〈某测量〉借用〈单位〉〈仪器〉完成，由〈姓名〉协助操作；
其余工作无他人协助。”——无他人协助也写明。〉

〈④ AI 使用披露：“本研究使用〈工具名及版本〉，用于〈环节：如文字润色 / 代码调试〉，
使用时间为〈时段〉，频率约〈次数〉；使用已事先获指导老师许可，相关聊天记录已上传
报名系统‘其他材料’栏。除上述外，研究设计、数据、推导与论文主体均未使用 AI。”
未使用则写：“本研究全过程未使用任何生成式 AI 工具。”〉

〈如有必要，可附参赛学生和指导老师的简历（官方模板致谢页第 1 条）。〉

\end{document}
